\documentclass[11pt]{article}
\usepackage[margin=1in]{geometry}
\usepackage{amsmath,amssymb}
\usepackage{xcolor}

\setlength{\parindent}{1.5em}

\newcommand{\approvalbox}{%
\begin{center}
\fbox{\begin{minipage}{0.85\textwidth}\centering
\textbf{Approval:} \quad $\square$ Student approved \qquad\qquad $\square$ Teaching staff approved
\end{minipage}}
\end{center}
\vspace{6pt}
}

\begin{document}

\begin{center}
{\Large Solutions --- Lab 1: Modeling and Identification of the Parrot AR.Drone}\\[6pt]
{\large Questions 2.1, 2.2, 2.3}
\end{center}

\vspace{18pt}


\vspace{18pt}

\section*{Setup and assumptions}

We use the notation of the handout: $\mathbf{f}_T = \begin{bmatrix} f_1 & f_2 & f_3 & f_4 \end{bmatrix}^T$ for the individual rotor thrusts, $n_i = c f_i$ for the individual reaction torques, $r$ for the arm length (distance from the center of mass to each rotor), and $\mathbf{f}_g = mg\mathbf{e}_3$ for gravity expressed in $\{I\}$.


\section*{2.1 -- Expressions for $\mathbf{f}$ and $\mathbf{n}$}

\approvalbox

\noindent\textbf{Translational force.} The total external force in $\{B\}$ is the sum of gravity (expressed in $\{I\}$, rotated into $\{B\}$) and the resultant of the four rotor thrusts (already aligned with $\{B\}$, all along $-z_B$):
\[
\mathbf{f} = R^T(\lambda)\,\mathbf{f}_g + N\mathbf{f}_T.
\]
Hence
\[
M = R^T(\lambda) \in SO(3), \qquad N = -\mathbf{e}_3 \begin{bmatrix} 1 & 1 & 1 & 1 \end{bmatrix} = \begin{bmatrix} 0 & 0 & 0 & 0 \\ 0 & 0 & 0 & 0 \\ -1 & -1 & -1 & -1 \end{bmatrix} \in \mathbb{R}^{3\times4}.
\]
($M$ depends on attitude because gravity is defined in $\{I\}$ and must be rotated into $\{B\}$; $N$ is constant because every thrust is already expressed in $\{B\}$ along the same axis.)

\noindent\textbf{Torque.} Each rotor $i$ produces (a) a moment from its thrust acting at moment arm $r$ about the two horizontal axes, and (b) a reaction torque $n_i = cf_i$ about $z_B$, alternating in sign with the spin direction. With the geometric convention above ($\mathbf{r}_i \times (-f_i\mathbf{e}_3)$ for each rotor position $\mathbf{r}_i$):
\[
n_x = r(f_2 - f_4), \qquad n_y = r(f_1 - f_3), \qquad n_z = c(f_1 - f_2 + f_3 - f_4),
\]
so that
\[
\mathbf{n} = P\mathbf{f}_T, \qquad P = \begin{bmatrix} 0 & r & 0 & -r \\ r & 0 & -r & 0 \\ c & -c & c & -c \end{bmatrix} \in \mathbb{R}^{3\times4}.
\]

\section*{2.2 -- Redefining $\{B\}$ with $x_B$ at $45^\circ$}

\approvalbox

Rotating the body axes by $45^\circ$ about $z_B$ (new $x_B'$ midway between the arms of rotors 1 and 2) does \emph{not} change $\mathbf{f}$: every thrust still acts along $-z_B$ (unchanged, since the rotation is about $z_B$ itself), and $R^T(\lambda)$ is still just the attitude rotation expressed with respect to whichever axes are called $\{B\}$ --- its functional form does not depend on this relabeling. Hence
\[
M,\ N \text{ unchanged.}
\]

$\mathbf{n} = P\mathbf{f}_T$ \emph{does} change, because the horizontal position of every rotor relative to the (now rotated) $x_B, y_B$ axes changes. The new $x_B'$ is midway between the old $+x_B$ (rotor 1's arm) and old $-y_B$ (rotor 2's arm), i.e. $x_B' = (x_B - y_B)/\sqrt{2}$, $y_B' = (x_B + y_B)/\sqrt{2}$. Expressing each rotor position in the new basis gives $\mathbf{r}_1' = \frac{r\sqrt{2}}{2}(1,1,0)$, $\mathbf{r}_2' = \frac{r\sqrt{2}}{2}(1,-1,0)$, $\mathbf{r}_3' = \frac{r\sqrt{2}}{2}(-1,-1,0)$, $\mathbf{r}_4' = \frac{r\sqrt{2}}{2}(-1,1,0)$ --- each rotor sits at $45^\circ$ from both axes, so each thrust now contributes to \emph{both} $n_x$ and $n_y$ (instead of purely one, as before). Using $\mathbf{r} \times (-f\mathbf{e}_3) = (-fr_y, fr_x, 0)$ for each rotor:
\[
n_x = \frac{r\sqrt{2}}{2}(f_2 + f_3 - f_1 - f_4), \qquad n_y = \frac{r\sqrt{2}}{2}(f_1 + f_2 - f_3 - f_4),
\]
while $n_z$ is unchanged (yaw reaction torque depends only on spin direction, not on the axes' in-plane orientation):
\[
n_z = c(f_1 - f_2 + f_3 - f_4).
\]
So
\[
P' = \begin{bmatrix} -\frac{r\sqrt{2}}{2} & \frac{r\sqrt{2}}{2} & \frac{r\sqrt{2}}{2} & -\frac{r\sqrt{2}}{2} \\[6pt] \frac{r\sqrt{2}}{2} & \frac{r\sqrt{2}}{2} & -\frac{r\sqrt{2}}{2} & -\frac{r\sqrt{2}}{2} \\[6pt] c & -c & c & -c \end{bmatrix} \neq P.
\]


\section*{2.3 -- Input transformation $\mathbf{u} = L\mathbf{f}_T$}

\approvalbox

Stacking $T = \sum_i f_i$ (from 2.1, first row all ones) on top of $\mathbf{n} = P\mathbf{f}_T$ (from 2.1) gives
\[
\mathbf{u} = \begin{bmatrix} T \\ \mathbf{n} \end{bmatrix} = L\mathbf{f}_T, \qquad L = \begin{bmatrix} 1 & 1 & 1 & 1 \\ 0 & r & 0 & -r \\ r & 0 & -r & 0 \\ c & -c & c & -c \end{bmatrix}.
\]

\noindent\textbf{Usefulness.} $L$ is a constant, invertible matrix ($\det L \neq 0$ for $r, c \neq 0$), so the map $\mathbf{f}_T \leftrightarrow \mathbf{u}$ is good. Its point is that the rigid-body dynamics (1)--(3) only ever depend on $\mathbf{f}_T$ through exactly these four combinations: $T$ enters the translational dynamics (scaled along $z_B$) and $\mathbf{n}$ enters the rotational dynamics directly. The individual $f_i$ are physically coupled (each one affects both $T$ and one or two components of $\mathbf{n}$); $\mathbf{u}$ decouples the problem into four independent virtual channels --- one for altitude/thrust and three for the body-frame moments (roll, pitch, yaw). This lets the height and attitude controllers be designed directly in terms of $T$ and $\mathbf{n}$ (as done throughout the rest of the guide), with the physical per-rotor commands recovered only at the very last step via $L^{-1}$ --- i.e., $L$ is exactly the control-allocation / mixer stage of a real flight controller. The only caveat is that not every $\mathbf{u}$ is achievable, since $\mathbf{f}_T$ must remain nonnegative and bounded.

\noindent\textbf{Inverse transformation.} Solving $L\mathbf{f}_T = \mathbf{u}$ directly (summing/subtracting pairs of rows):
\[
f_1 = \frac{T}{4} + \frac{n_y}{2r} + \frac{n_z}{4c}, \qquad f_2 = \frac{T}{4} + \frac{n_x}{2r} - \frac{n_z}{4c},
\]
\[
f_3 = \frac{T}{4} - \frac{n_y}{2r} + \frac{n_z}{4c}, \qquad f_4 = \frac{T}{4} - \frac{n_x}{2r} - \frac{n_z}{4c},
\]
i.e. $\mathbf{f}_T = L^{-1}\mathbf{u}$ with
\[
L^{-1} = \begin{bmatrix} 1/4 & 0 & 1/(2r) & 1/(4c) \\ 1/4 & 1/(2r) & 0 & -1/(4c) \\ 1/4 & 0 & -1/(2r) & 1/(4c) \\ 1/4 & -1/(2r) & 0 & -1/(4c) \end{bmatrix}.
\]

Each rotor receives a quarter of the total thrust plus/minus the roll, pitch and yaw contributions needed to produce the requested $\mathbf{n}$ --- the familiar quadrotor mixer formula.

\end{document}
